\chapter{Introducción específica} % Main chapter title

\label{Chapter2}

El desarrollo de una solución basada en el procesamiento de lenguaje natural exige un acoplamiento profundo y eficiente con el hardware del dispositivo móvil. Este apartado detalla el ecosistema tecnológico y de software seleccionado para el desarrollo del prototipo.

%----------------------------------------------------------------------------------------
%	SECTION 1
%----------------------------------------------------------------------------------------

\section{Marco metodológico}

Para la concepción, desarrollo y evaluación del sistema de transacciones financieras por voz, se adoptó como marco de trabajo la metodología CRISP-DM (\textit{Cross-Industry Standard Process for Data Mining}) \citep{wirth2000crisp}. Si bien CRISP-DM fue concebida originalmente para proyectos de minería de datos, su estructura iterativa y centrada en el ciclo de vida de los datos la convierte en el estándar de facto para la ingeniería de proyectos de inteligencia artificial y machine learning. La elección de esta metodología responde a la necesidad de contar con un enfoque estructurado que permita alinear los objetivos técnicos del procesamiento de lenguaje natural con los requisitos críticos y las restricciones de seguridad propios de una aplicación de banca móvil.
El ciclo de vida se ha adaptado específicamente para abordar el desarrollo de un pipeline NLU (\textit{Natural Language Understanding}) \citep{chen2019bert} de ejecución local dentro del ecosistema iOS. 


\section{Herramientas utilizadas}

Con el objetivo de cumplir con los requerimientos de procesamiento local (\textit{Edge AI}) \citep{shi2016edge} y garantizar la privacidad por diseño (\textit{Privacy by Design}) \citep{cavoukian2009privacy} en el contexto de una aplicación financiera, se adoptó de manera integral el entorno nativo de Apple. A continuación, se enumeran las herramientas integradas en la arquitectura del sistema.

\begin{itemize}

    \item Swift \citep{apple2026swiftlanguage}: es un lenguaje de programación compilado, multiparadigma y de tipado estático fuerte, desarrollado por Apple para la creación de software en sus sistemas operativos. Su adopción se justifica por su alto rendimiento computacional y su seguridad en el manejo de memoria mediante el Conteo Automático de Referencias (ARC). Estas características son esenciales para manipular flujos de datos en tiempo real (como el audio) y procesar estructuras de datos complejas sin incurrir en fugas de memoria o vulnerabilidades durante la ejecución transaccional.

    \item Xcode \citep{apple2026xcode}: es el Entorno de Desarrollo Integrado (IDE) oficial y nativo provisto por Apple, que incluye los compiladores, depuradores y simuladores necesarios para el ciclo de vida del software. Se utilizó como la herramienta central para la codificación, compilación y empaquetado del SDK. Su uso es mandatorio para desplegar la prueba de concepto en dispositivos físicos y para utilizar la herramienta \textit{Instruments}, esta permitió perfilar y auditar el consumo de memoria y CPU de los modelos de Inteligencia Artificial durante la inferencia local.

    \item AVFoundation \citep{apple2026avfoundation}: es el \textit{framework} nativo de Apple que proporciona una interfaz de programación de bajo nivel para interactuar con el hardware audiovisual del dispositivo. Su uso fue indispensable para orquestar la captura segura del comando de voz. Específicamente, se utilizó la clase \texttt{AVAudioSession} para gestionar los permisos de privacidad del micrófono a nivel del sistema operativo, y la clase \texttt{AVAudioEngine} para interceptar la onda sonora cruda y aislar el flujo acústico sin depender de librerías de terceros.

    \item Speech Framework \citep{apple2026speech}: es una Interfaz de Programación de Aplicaciones (API) de Apple dedicada exclusivamente al Reconocimiento Automático de Habla (ASR), capaz de transcribir señales de audio a cadenas de texto. Se justificó su uso frente a otras alternativas comerciales debido a su capacidad de operar en modo \textit{offline}. Al forzar la propiedad \lstinline|requiresOnDevice-Recognition|, esta herramienta garantizó que la voz del usuario nunca fuera transmitida a servidores externos, mitigando los vectores de ataque y cumpliendo con el estándar de secreto bancario requerido por el proyecto.

    \item Core ML \citep{apple2026coreml}: es el \textit{framework} central de \textit{Machine Learning} de Apple, diseñado para integrar, optimizar y ejecutar modelos predictivos directamente en el sistema operativo iOS. Constituye el motor analítico del proyecto. Se seleccionó porque delega automáticamente las cargas matemáticas pesadas (como la clasificación de intenciones y la extracción de entidades) al \textit{Apple Neural Engine} (ANE). Esto garantiza que los algoritmos de Procesamiento de Lenguaje Natural (PLN) se ejecuten con una latencia mínima y un impacto energético marginal, habilitando el paradigma de inferencia 100\% \textit{On-Device}.

    \item Coremltools \citep{apple2026coremltools}: es un paquete de código abierto basado en Python, desarrollado por Apple, que unifica la conversión de modelos entrenados en librerías estándar (como PyTorch o TensorFlow) al formato propietario \linebreak \lstinline|.mlpackage|. Su aplicación metodológica fue crítica para viabilizar el despliegue de las redes neuronales en el entorno móvil. Se utilizó para aplicar técnicas de cuantización a los pesos sinápticos de los modelos, reduciendo drásticamente su tamaño de almacenamiento y la huella en la memoria RAM del dispositivo, sin comprometer significativamente la exactitud predictiva del sistema.

\end{itemize}

\section{Requerimientos y modelos de IA utilizados}

El núcleo algorítmico del prototipo aborda la comprensión del comando de voz como un problema predictivo dual de machine learning: la clasificación de intenciones y el etiquetado de entidades transaccionales.

Por un lado, la clasificación multiclase a nivel de oración tiene como objetivo inferir la acción principal del usuario (por ejemplo, iniciar una transferencia o consultar un saldo). Matemáticamente, la red neuronal proyecta la secuencia de texto vectorizada hacia un espacio discreto de clases predefinidas, empleando una función de activación \textit{softmax} para seleccionar, de manera excluyente, la intención que maximiza la probabilidad estadística.

Por otro lado, la conformación del \textit{payload} financiero exige resolver un problema de etiquetado de secuencias a nivel de \textit{token}, enmarcado en el Reconocimiento de Entidades Nombradas (NER)\citep{lample2016neural}. En lugar de emitir un único juicio para toda la oración, el modelo asigna una clase categórica a cada palabra utilizando el esquema de anotación IOB (\textit{Inside, Outside, Beginning}). Esta técnica permite al sistema demarcar los límites exactos de variables críticas, como el monto numérico, la divisa y el destinatario, capturando las dependencias semánticas del lenguaje para extraer los parámetros exactos de la operación.



\section{Exploración teórica de arquitecturas evaluadas}

Para abordar este desafío predictivo dual y que se respeten las restricciones de memoria y procesamiento de los dispositivos móviles, se llevó a cabo una evaluación teórica de diversas arquitecturas algorítmicas, que analiza el compromiso (\textit{trade-off}) entre precisión semántica, latencia de inferencia y consumo de recursos.

En primer lugar, se analizó el estado del arte representado por las arquitecturas basadas en Transformers. Este paradigma, fundamentado en el mecanismo de autoatención, prescinde de la recurrencia secuencial y permite capturar dependencias a largo plazo y procesa el contexto de manera bidireccional y paralelizable. Dentro de esta familia, se evaluó la viabilidad de incorporar modelos del lenguaje preentrenados adaptados al español, como BETO (la variante hispanohablante de BERT) \citep{canete2020spanish}. Si bien BETO ofrece un rendimiento superior en la comprensión de las sutilezas léxicas y morfológicas del lenguaje coloquial rioplatense, su implementación nativa en el edge presenta obstáculos significativos. La alta dimensionalidad de sus parámetros (típicamente superior a 110 millones) exige la aplicación de técnicas agresivas de compresión, poda (\textit{pruning}) y destilación de conocimiento (\textit{knowledge distillation}) para transformarlo en un artefacto viable para el entorno iOS sin agotar la memoria RAM del dispositivo.

Como alternativa de menor huella computacional, se exploró la familia de las redes neuronales recurrentes (RNN) \citep{elman1990finding}. A diferencia de los transformers, las RNN procesan los tokens de manera estrictamente secuencial, actualiza un estado oculto que actúa como memoria temporal del contexto lingüístico. Dado que las arquitecturas RNN tradicionales son susceptibles al problema del desvanecimiento del gradiente (\textit{vanishing gradient}) al procesar secuencias largas, la investigación se centró en sus variantes con mecanismos de compuerta (\textit{gating}): las redes de memoria a corto y largo plazo (LSTM) \citep{hochreiter1997long} y las unidades recurrentes con compuerta (GRU) \citep{cho2014learning}. Estas topologías regulan matemáticamente el flujo de información, dice qué características del contexto previo deben retenerse y cuáles descartarse. Históricamente, las arquitecturas basadas en LSTM bidireccionales (BiLSTM) \citep{huang2015bidirectional} han demostrado una alta eficacia para el etiquetado de secuencias (NER) \citep{lample2016neural}, logra una precisión competitiva con un costo computacional y tamaño de modelo órdenes de magnitud inferiores a los de un modelo fundacional, lo que las hace inherentemente más compatibles con inferencias On-Device \citep{dhar2021survey}.

Finalmente, se evaluó la pertinencia de utilizar los frameworks de comprensión de lenguaje natural nativos de Apple, en particular la API NaturalLanguage y el ecosistema de Create ML \citep{apple2025naturallanguage}. Apple proporciona clases altamente especializadas, como NLModel y NLTagger, diseñadas para compilar clasificadores de texto y extractores de entidades a partir de \textit{datasets} etiquetados. La ventaja superlativa de este enfoque reside en su integración opaca y de bajo nivel con el silicio del dispositivo. Los modelos generados bajo este estándar están nativamente optimizados para aprovechar la aceleración por hardware del procesamiento neural de Apple, garantiza tiempos de latencia del orden de los milisegundos y un impacto insignificante en la autonomía de la batería. Aunque la utilización de frameworks cerrados sacrifica el grado de control arquitectónico y la hiperparametrización profunda que ofrecen librerías como PyTorch o TensorFlow, su robustez operativa, la garantía de privacidad al aislar el procesamiento de la red, y la facilidad de integración mediante Core ML, posicionan a estas herramientas nativas como candidatos estratégicos para validar la viabilidad técnica de una prueba de concepto en el sector financiero.

A continuación se detallan las métricas que respaldan lo mencionado en el presente capítulo:

% =============================================================================
% tablas_metricas.tex
%
% Tabla LaTeX para la sección de evaluación de modelos de la tesis
% "Sistema de transacciones financieras por voz en dispositivos móviles iOS"
%
% Ing. Alexis Geraldine Barniquez Piñero - FIUBA
% =============================================================================


% =============================================================================
% TABLA 1: COMPARATIVA DE ARQUITECTURAS (Trade-off)
% =============================================================================
% Esta tabla respalda la discusión teórica de la sección 2.x sobre la
% evaluación de arquitecturas: Transformers vs RNN vs Apple Native.

\begin{table}[htbp]
\centering
\resizebox{\textwidth}{!}{%
\begin{tabular}{@{}lcccc@{}}
\toprule
\textbf{Criterio}
  & \textbf{BETO}
  & \textbf{BiLSTM}
  & \textbf{Create ML}
  & \textbf{Create ML} \\
  & \textbf{(Transformer)}
  & \textbf{(RNN)}
  & \textbf{(MaxEnt)}
  & \textbf{(CRF)} \\
\midrule
Tamaño del modelo (MB)
  & $\sim$440
  & $\sim$5--15
  & $<$1
  & $<$1 \\
Parámetros (millones)
  & $\sim$110
  & $\sim$0.5--2
  & N/A
  & N/A \\
Latencia de inferencia (ms)
  & $>$200
  & 20--50
  & $<$10
  & $<$10 \\
Memoria RAM en runtime (MB)
  & $>$300
  & 15--40
  & $<$5
  & $<$5 \\
Requiere conversión Core ML
  & Sí
  & Sí
  & No
  & No \\
Aceleración Neural Engine
  & Parcial
  & Parcial
  & Nativa
  & Nativa \\
Entrenamiento sin GPU
  & No viable
  & Posible
  & Sí (CPU)
  & Sí (CPU) \\
Privacidad \textit{on-device}
  & Sí
  & Sí
  & Sí
  & Sí \\
\bottomrule
\end{tabular}%
}
\caption{Comparativa de arquitecturas evaluadas para NLU \textit{on-device} en iOS.
         Los valores de BETO y BiLSTM corresponden a referencias bibliográficas;
         los de Create ML (MaxEnt/CRF) son resultados experimentales propios.}
\label{tab:comparativa-arquitecturas}
\end{table}


